\documentclass[11pt]{amsart}
\usepackage[localtheorems]{raeez-math-template}

\newtheorem{theorem}{Theorem}[section]
\newtheorem{proposition}[theorem]{Proposition}
\newtheorem{lemma}[theorem]{Lemma}
\newtheorem{corollary}[theorem]{Corollary}
\theoremstyle{definition}
\newtheorem{definition}[theorem]{Definition}
\newtheorem{hypothesis}[theorem]{Hypothesis}
\newtheorem{example}[theorem]{Example}
\theoremstyle{remark}
\newtheorem{remark}[theorem]{Remark}

\providecommand{\En}{\mathsf{E}_n}
\providecommand{\Eone}{\mathsf{E}_1}
\providecommand{\Etwo}{\mathsf{E}_2}
\providecommand{\Ethree}{\mathsf{E}_3}
\providecommand{\cC}{\mathcal{C}}
\providecommand{\cO}{\mathcal{O}}
\providecommand{\cE}{\mathcal{E}}
\providecommand{\cF}{\mathcal{F}}
\providecommand{\cZ}{\mathcal{Z}}
\providecommand{\CC}{\mathbb{C}}
\providecommand{\RR}{\mathbb{R}}
\providecommand{\ZZ}{\mathbb{Z}}
\providecommand{\HH}{\mathrm{HH}}
\providecommand{\HKR}{\mathrm{HKR}}
\providecommand{\KT}{\mathrm{KT}}
\providecommand{\BV}{\mathrm{BV}}
\providecommand{\PV}{\mathrm{PV}}
\providecommand{\Tpoly}{T_{\mathrm{poly}}}
\providecommand{\Dpoly}{D_{\mathrm{poly}}}
\providecommand{\Perf}{\mathrm{Perf}}
\providecommand{\Rep}{\mathrm{Rep}}
\providecommand{\Conf}{\mathrm{Conf}}
\providecommand{\FM}{\mathrm{FM}}
\providecommand{\GC}{\mathrm{GC}}
\providecommand{\hCS}{\mathrm{hCS}}
\providecommand{\Obs}{\mathrm{Obs}}
\providecommand{\QME}{\mathrm{QME}}
\providecommand{\MC}{\mathrm{MC}}
\providecommand{\dd}{\mathrm{d}}
\providecommand{\id}{\mathrm{id}}
\providecommand{\Sym}{\mathrm{Sym}}
\providecommand{\End}{\mathrm{End}}
\providecommand{\Schouten}{\lbrack\!\lbrack-,-\rbrack\!\rbrack}
\providecommand{\GraphEdges}{E}
\providecommand{\vol}{\mathrm{vol}}

\title{Koszul--Tate resolutions, formality, and the $d=3$ operadic
boundary}
\author{Raeez Lorgat}
\date{2026-04-22}

\begin{document}
\maketitle

\begin{abstract}
The Koszul--Tate resolution, the Hochschild--Kostant--Rosenberg map,
Kontsevich--Tamarkin formality, operadic Koszul duality, and the
CY-to-chiral comparison solve different problems.  The exact interface
at dimension three is as follows.  The local configuration-space
calculation gives the degree \(1-3=-2\) bracket in the homology operad
\(e_3=H_\bullet(\mathsf{D}_3)\).  It supplies a chain-level
formality input only under the usual graph-integral hypotheses.  The
specialised CY-to-chiral output at \(d=3\) is an \(\Eone\)-chiral
algebra; braided or higher little-disks structure belongs to centre,
double, Hochschild, or module layers after their own hypotheses are
fixed.
\end{abstract}

\section{Five Constructions}
\label{sec:five-constructions}

Let \(k\) be a field of characteristic zero.  The following objects are
kept separate throughout.

\begin{definition}[Koszul--Tate resolution]
\label{def:kt-resolution}
Let \(R\) be a commutative \(k\)-algebra and let
\(f=(f_1,\ldots,f_r)\subset R\).  The elementary Koszul--Tate complex is
the non-positive cdga
\[
  \KT(R;f)=\bigl(R[\eta_1,\ldots,\eta_r],\; \dd \eta_a=f_a,\;
  |\eta_a|=-1\bigr).
\]
If \(f_1,\ldots,f_r\) is a regular sequence, then
\[
  H^0\KT(R;f)\cong R/(f_1,\ldots,f_r),\qquad
  H^i\KT(R;f)=0\quad (i<0).
\]
For a non-regular system one adds higher generators to kill the
negative homology.  The resulting cdga resolves the derived zero locus
of the equations.  Formality requires separate homotopy-transfer data.
\end{definition}

\begin{definition}[HKR]
\label{def:hkr}
For \(X\) smooth over \(k\), Hochschild--Kostant--Rosenberg gives the
graded vector-space identification
\[
  \HKR_X\colon
  \HH^i(\Perf(X))
  \xrightarrow{\;\sim\;}
  \bigoplus_{p+q=i} H^q\bigl(X,\wedge^p T_X\bigr).
\]
On affine smooth schemes this is induced by sending a polyvector field
to the antisymmetrised Hochschild cochain.  For smooth proper schemes
the Todd-class correction belongs to trace, Mukai-pairing, Chern
character, and global module compatibilities.  HKR is a cohomology
identification; it does not provide a chain-level \(\En\)-formality
quasi-isomorphism.
\end{definition}

\begin{definition}[Formality quasi-isomorphism]
\label{def:formality}
A formality statement is a quasi-isomorphism between a chain object and
its homology object, equipped with the relevant operadic structure.
For little disks this takes the form
\[
  C_\bullet(\mathsf{D}_n;k)\xrightarrow{\;\simeq\;}
  H_\bullet(\mathsf{D}_n;k)=e_n
\]
over characteristic zero, with graph-integral models over \(\RR\).
For deformation quantisation, Kontsevich's local map
\[
  U\colon \Tpoly(\RR^m)\longrightarrow \Dpoly(\RR^m)
\]
is an \(L_\infty\) quasi-isomorphism built from admissible graphs.
Tamarkin's theorem upgrades the Hochschild side to Gerstenhaber
formality using an associator.  These are formality theorems, not
Koszul--Tate resolutions and not HKR itself.
\end{definition}

\begin{definition}[Operadic Koszul duality and BV]
\label{def:operadic-bv}
Operadic Koszul duality for little disks says
\[
  \En^{!}\simeq \En\{n\},
\]
up to the standard operadic suspension convention.  This accounts for
the bracket in \(e_n\) having cohomological degree \(1-n\).
The BV formalism is a separate package: a shifted symplectic space of
fields, a classical differential \(Q\), an interaction \(I\), an odd
Poisson bracket \(\{-,-\}\), and, after gauge fixing, a BV Laplacian
\(\Delta\).  Perturbative quantisation requires the quantum master
equation
\[
  QI+\frac12\{I,I\}+\hbar\Delta I=0.
\]
The Koszul--Tate complex may appear inside the classical BV complex as
the resolution of equations of motion.  Operadic Koszul duality explains
bar-cobar and shifted-bracket conventions.  The two mechanisms have
different inputs.
\end{definition}

\begin{definition}[CY-to-chiral comparison at \(d=3\)]
\label{def:cy-to-chiral-d3}
Let \(\Phi_3\) denote the \(d=3\) specialisation of the CY-to-chiral
construction on its framed, strictified domain.  Its output is an
\(\Eone\)-chiral algebra
\[
  A=\Phi_3(\cC).
\]
Braided \(\Etwo\)-structure appears on layers such as
\[
  \cZ\bigl(\Rep^{\Eone}(A)\bigr),
\]
and higher Hochschild or factorisation structures appear only after the
corresponding centre, double, module, or higher-Deligne hypotheses have
been supplied.  The local \(\Ethree\)-operadic calculation below is an
input to those auxiliary layers.  The native operadic level of \(A\)
remains \(\Eone\).
\end{definition}

\section{Configuration-Space Calculation at \(n=3\)}
\label{sec:configuration}

The calculation needed in dimension three is elementary and useful.
It proves a local degree statement for \(e_3\) and a low-order
vanishing statement for graph weights with no external ground
vertices.

\begin{definition}[Fulton--MacPherson space and angle form]
\label{def:fm-angle}
For \(N\geq 2\), let
\[
  \Conf_N(\RR^3)=\{(x_1,\ldots,x_N)\in(\RR^3)^N:x_i\neq x_j
  \text{ for } i\neq j\}.
\]
After quotienting by translations and positive dilations, its
Fulton--MacPherson compactification is denoted \(\FM_3(N)\).  Its real
dimension is
\[
  \dim \FM_3(N)=3N-4.
\]
For an ordered pair \(i\neq j\), the map
\[
  \varphi_{ij}\colon \Conf_N(\RR^3)\longrightarrow S^2,\qquad
  \varphi_{ij}(x)=\frac{x_j-x_i}{|x_j-x_i|}
\]
pulls the unit-volume form on \(S^2\) back to a closed \(2\)-form
\[
  \omega_{ij}=\varphi_{ij}^{*}\vol_{S^2}.
\]
The form extends smoothly to \(\FM_3(N)\).
\end{definition}

\begin{definition}[Graph weight with no ground vertices]
\label{def:graph-weight}
For a directed admissible graph \(\Gamma\) on \(N\) labelled vertices,
with no self-loops and no repeated directed edges, set
\[
  w_\Gamma =
  \int_{\FM_3(N)} \bigwedge_{e=(i\to j)\in \GraphEdges(\Gamma)}\omega_{ij}
\]
when the wedge has top degree, and set \(w_\Gamma=0\) otherwise.  A
normalising power of \(2\pi\) may be inserted by convention; the
top-degree criterion is unchanged.
\end{definition}

\begin{proposition}[Degree count]
\label{prop:degree-count}
If \(w_\Gamma\) is a top-degree integral on \(\FM_3(N)\), then
\[
  2|\GraphEdges(\Gamma)|=3N-4.
\]
Consequently, graphs with \(N=2\) require one edge, and graphs with
\(N=3\) give zero weights in the no-ground-vertex sector.
\end{proposition}

\begin{proof}
Each edge contributes the \(2\)-form \(\omega_{ij}\).  The quotient
configuration space has dimension \(3N-4\).  Equality of form degree and
manifold dimension is therefore the displayed equation.  For \(N=2\)
one obtains \(|\GraphEdges(\Gamma)|=1\).  For \(N=3\) one obtains
\(|\GraphEdges(\Gamma)|=5/2\), which is impossible for a graph.
\end{proof}

\begin{proposition}[The binary operation]
\label{prop:binary-operation}
The one-edge graph at \(N=2\) integrates over
\(\FM_3(2)\cong S^2\), and
\[
  \int_{\FM_3(2)}\omega_{12}=1
\]
with the unit-volume normalisation.  Its homology class is the bracket
generator in \(e_3\), hence has cohomological degree
\[
  1-3=-2.
\]
\end{proposition}

\begin{proof}
After translation and dilation, a pair of distinct points in \(\RR^3\)
is determined by the direction from the first point to the second,
which is \(S^2\).  The angle form is the unit-volume form on this
sphere.  Cohen's description of \(H_\bullet(\mathsf{D}_n)\) gives an
\((n-1)\)-homological bracket, equivalently a cohomological bracket of
degree \(1-n\).  At \(n=3\) this degree is \(-2\).
\end{proof}

\begin{remark}[Relation to the Schouten bracket]
\label{rem:schouten-degree}
HKR identifies Hochschild cohomology with polyvector cohomology.  The
ordinary Schouten--Nijenhuis bracket on polyvectors has its usual
Gerstenhaber degree.  Comparing it with the \(e_3\)-bracket requires the
explicit suspension dictated by the \(e_3\) convention.  Without that
suspension, the formula is the Schouten bracket as a bilinear
differential operator; the degree statement belongs to the shifted
\(e_3\)-algebra.
\end{remark}

\section{Local Formality and Its Boundary}
\label{sec:local-formality}

\begin{theorem}[Local little-disks formality, dimension three]
\label{thm:local-e3-formality}
Over characteristic zero, the chain operad of little \(3\)-disks is
formal:
\[
  C_\bullet(\mathsf{D}_3;k)\xrightarrow{\;\simeq\;}e_3.
\]
Over \(\RR\), graph-integral models realise this quasi-isomorphism by
closed angle forms on Fulton--MacPherson compactifications.  The
boundary-stratum identities are Stokes identities for the closed forms
\(\omega_{ij}\).
\end{theorem}

\begin{proof}
Kontsevich's configuration-space proof constructs closed forms on
\(\FM_3(N)\) representing the generators and relations of \(e_3\).
Since \(\dd\omega_{ij}=0\), Stokes' theorem gives
\[
  0=\int_{\FM_3(N)}\dd\Bigl(\bigwedge_e\omega_e\Bigr)
   =\int_{\partial\FM_3(N)}\bigwedge_e\omega_e.
\]
Each boundary stratum is an operadic composition of smaller
configuration spaces.  The resulting identities are exactly the
quadratic and higher homotopy identities of the graph model.  The
quasi-isomorphism follows from Kontsevich's formality theorem for
little disks.
\end{proof}

\begin{hypothesis}[What global use requires]
\label{hyp:global-formality}
A global use of Theorem~\ref{thm:local-e3-formality} on a smooth
Calabi--Yau threefold \(X\) requires additional data:
\begin{enumerate}
\item a sheaf or factorisation-algebra model whose local charts are the
formal disks used by the graph construction;
\item descent maps on overlaps, with the Atiyah and Todd corrections
specified;
\item a completion in which all infinite graph sums converge or are
interpreted perturbatively;
\item a compatibility statement between the chain-level model and the
chosen \((\infty,1)\)-categorical Hochschild or centre construction.
\end{enumerate}
The vanishing \(c_1(X)=0\) is one input to these compatibilities.  It
does not erase every Todd or Atiyah correction by itself.
\end{hypothesis}

\begin{corollary}[Local consequences of the low-order calculation]
\label{cor:local-consequences}
The graph calculation in Section~\ref{sec:configuration} proves the
following local statements.
\begin{enumerate}
\item The no-ground-vertex sector has no cubic weight at \(N=3\).
\item The binary bracket in \(e_3\) has cohomological degree \(-2\).
\item HKR can identify the underlying cohomology class with the
Schouten operation after the chosen suspension is fixed.
\end{enumerate}
It does not construct a global \(\Ethree\)-algebra on
\(\Phi_3(\cC)\).
\end{corollary}

\section{Holomorphic Chern--Simons}
\label{sec:hcs}

Let \(X\) be a Calabi--Yau threefold with holomorphic volume form
\(\Omega_X\), and let \(\mathfrak g\) be a finite-dimensional Lie
algebra equipped with an invariant pairing.  The classical hCS fields
on \(X\) are
\[
  \cE_X=\Omega^{0,\bullet}(X,\mathfrak g)[1],
\]
with action
\[
  S(A)=\int_X \Omega_X\wedge
  \left(\frac12\langle A,\bar\partial A\rangle
  +\frac16\langle A,[A,A]\rangle\right).
\]
The linearised Koszul--Tate part resolves the equations of motion near
a chosen classical solution.  The BV bracket and the cubic interaction
encode gauge symmetry and deformation theory.

\begin{hypothesis}[Perturbative quantisation of compact hCS]
\label{hyp:hcs-qme}
A compact hCS comparison with graph formality requires:
\begin{enumerate}
\item an orientation datum from \(\Omega_X\) and the invariant pairing;
\item a gauge fixing or framing and a heat-kernel regularisation;
\item a completed space of local functionals;
\item vanishing or cancellation of the local anomaly class;
\item a solution of the quantum master equation
\[
  QI+\frac12\{I,I\}+\hbar\Delta I=0.
\]
\end{enumerate}
Under these hypotheses, perturbative Feynman weights are legitimate
configuration-space integrals.  Matching them with the
Kontsevich--Tamarkin weights is then a theorem of the chosen local
model.  HKR and the Koszul--Tate resolution supply only the indicated
input complexes.
\end{hypothesis}

\begin{remark}[Flat model]
\label{rem:flat-hcs}
On \(\CC^3\) with its standard holomorphic volume form, Costello--Li
construct perturbative hCS/BCOV quantisations by BV renormalisation.
The propagator has the expected diagonal singularity and gives
configuration-space weights after regularisation.  Extending a flat
graph equality to compact \(X\) asks for the five hypotheses in
Hypothesis~\ref{hyp:hcs-qme}.
\end{remark}

\section{CY-to-Chiral Consequence}
\label{sec:cy-to-chiral-consequence}

\begin{theorem}[Dimension-three separation]
\label{thm:d3-separation}
Let \(\cC\) be a Calabi--Yau \(3\)-category in the domain of the
specialised functor \(\Phi_3\).  The object
\[
  A=\Phi_3(\cC)
\]
is an \(\Eone\)-chiral algebra.  The local \(\Ethree\)-formality
calculation may be used in the construction or verification of higher
Hochschild, centre, double, or module structures associated to \(A\),
provided the corresponding descent, completion, and obstruction
hypotheses are supplied.  It does not change the native operadic level
of \(A\).
\end{theorem}

\begin{proof}
The \(d=3\) specialisation integrates or restricts along two complex
directions and leaves an ordered one-dimensional chiral output.  That
is the \(\Eone\) layer.  Dunn additivity and higher Deligne-type
theorems can create \(\Etwo\) or higher structures on centres and
Hochschild-type constructions, since these are additional algebraic
objects built from \(A\).  The little \(3\)-disks formality theorem
controls the local chain model for such higher structures.  The output
of \(\Phi_3\) remains \(\Eone\)-chiral.
\end{proof}

\begin{corollary}[No invariant is involved]
\label{cor:no-invariants}
The arguments in this note do not compute
\(\kappa_{\mathrm{ch}}\), \(\kappa_{\mathrm{cat}}\),
\(\kappa_{\mathrm{BKM}}=c_N(0)/2\), or
\(\kappa_{\mathrm{fiber}}\).  Those four invariants enter separate
CY-landscape and Borcherds-weight calculations.
\end{corollary}

\section{Citations}
\label{sec:citations}

\begin{itemize}
\item Hochschild, Kostant, Rosenberg.  Differential forms on regular
affine algebras, \emph{Transactions of the AMS} 102, 1962.
\item Kontsevich.  Deformation quantization of Poisson manifolds, I,
arXiv:q-alg/9709040.
\item Kontsevich.  Operads and motives in deformation quantization,
arXiv:math/9904055.
\item Tamarkin.  Another proof of M. Kontsevich formality theorem,
arXiv:math/9803025.
\item Cohen.  The homology of \(C_{n+1}\)-spaces, \(n\geq 0\), in
\emph{The Homology of Iterated Loop Spaces}, 1976.
\item Fresse.  Koszul duality of \(E_n\)-operads, arXiv:0809.1928.
\item Calaque and Van den Bergh.  Hochschild cohomology and Atiyah
classes, arXiv:0708.2725.
\item Costello and Gwilliam.  \emph{Factorization Algebras in Quantum
Field Theory}, Vols. 1 and 2.
\item Costello and Li.  Quantization of open-closed BCOV theory, I,
arXiv:1505.06703.
\end{itemize}

\end{document}
